% The double-tail latch (Schinkel) as flow/comp/subckt.py builds it: cross-coupled PMOS (latch_on)
% on a clocked PMOS tail (twice latch_on), cross-coupled NMOS (latch_init) to ground, and NMOS
% (rst) on the first stage's drains P and Q that hold the latch nodes low during reset.
% The preamble of every isometric drawing: standalone where installed (as docs/images/*.tex),
% else article with the page cropped to the drawing. Then Latin Modern sans and the blocks.
\IfFileExists{standalone.cls}{\documentclass[tikz,border=2mm]{standalone}\newenvironment{drawing}{}{}}{%
  \documentclass{article}\usepackage{tikz}\pagestyle{empty}\newsavebox\drawingbox
  \newenvironment{drawing}{
  \begin{lrbox}{\drawingbox}}{
    \end{lrbox}%
    \pagewidth=\dimexpr\wd\drawingbox+4mm\relax\pageheight=\dimexpr\ht\drawingbox+\dp\drawingbox+4mm\relax
\hoffset=-1in\voffset=-1in\shipout\vbox{\kern2mm\hbox{\kern2mm\usebox\drawingbox}}}}
\usepackage[T1]{fontenc}
\usepackage{lmodern}
\renewcommand{\familydefault}{\sfdefault}
\usepackage{thesis-drawings}

\usepackage{thesis-circuits}
\begin{document}
\begin{drawing}
  \begin{tikzpicture}[circuit]
    \draw (0,4.6) node[pmos](mo){};
    \draw (mo.S) node[vcc]{};
    \draw (mo.G) node[left, sig]{$\overline{\mathrm{CLK}}$};
    \noteR{mo}{$2\times$on}
    \draw (-1.6,3.83) -- (1.6,3.83);
    \fill (mo.D) circle (1.6pt);
    \draw (-1.6,3.06) node[pmos, xscale=-1](ml){};
    \draw (1.6,3.06) node[pmos](mr){};
    \draw (-1.6,0.96) node[nmos, xscale=-1](nl){};
    \draw (1.6,0.96) node[nmos](nr){};
    \draw (nl.S) node[ground]{} (nr.S) node[ground]{};
    \draw (ml.D) -- (nl.D) (mr.D) -- (nr.D);
    \draw (ml.G) -- (nl.G) (mr.G) -- (nr.G);
    \coordinate (X) at ($(nl.D)!0.7!(ml.D)$);
    \coordinate (Y) at ($(nr.D)!0.3!(mr.D)$);
    \draw (X) -- (X -| nl.G) node[jump crossing]{} -- (X -| nr.G);
    \draw (Y) -- (Y -| nr.G) node[jump crossing]{} -- (Y -| nl.G);
    \fill (X) circle (1.6pt) (Y) circle (1.6pt) (X -| nr.G) circle (1.6pt) (Y -| nl.G) circle (1.6pt);
    \draw ($(X) - (1.4,0.77)$) node[nmos](il){};
    \draw ($(Y) + (1.4,-0.77)$) node[nmos, xscale=-1](ir){};
    \draw (il.S) node[ground]{} (ir.S) node[ground]{};
    \draw (il.D) -- (X) (ir.D) -- (Y);
    \draw (il.G) node[left, sig]{$P$};
    \draw (ir.G) node[right, sig]{$Q$};
    \noteL{ml}{on}
    \noteL{nl}{init}
    \noteR{il}{rst}
    \draw (X) node[above left, sig]{$X$};
    \draw (Y) node[below right, sig]{$Y$};
  \end{tikzpicture}
\end{drawing}
\end{document}
